(1) 绝热理想气体满足 $TV^{\gamma-1}=$ 常数，因此压缩段 $T_2=T_1r^{\gamma-1}$。膨胀段体积比 $V_4/V_3=r$，所以 $T_4=T_3/r^{\gamma-1}$。两条定容线意味着 $V_2=V_3$、$V_1=V_4$；这些等式是奥托循环而非狄塞尔循环的关键。
(2) 定容热量为 $q_{\rm in}=c_v(T_3-T_2)$、$q_{\rm out}=c_v(T_4-T_1)$。利用 $T_4-T_1=(T_3-T_2)/r^{\gamma-1}$，热效率 $\eta=1-q_{\rm out}/q_{\rm in}=1-r^{1-\gamma}$，与热输入绝对量无关，前提为定热容与两绝热过程可逆。
(3) 第一定律对一个闭合循环有 $\Delta u=0$，故比净功 $w_{\rm net}=q_{\rm in}-q_{\rm out}=c_v(T_3-T_2)(1-r^{1-\gamma})$。运行时应有 $T_3>T_2$，对应正吸热、正净功；若 $T_3=T_2$，循环退化为纯压缩膨胀往返而不输出净功。
(4) 当 $r\to1^+$，$\eta=1-r^{1-\gamma}\to0$，热机失去有效压缩膨胀环节。理想式随 $r$ 提高而提升，但真实发动机受爆震、自燃、热应力、摩擦、换热损失和可变热容限制，不能将 $\eta\to1$ 当作可实现的无代价极限。
